%% ****** Start of file apsguide4-2.tex ****** %
%%
%%   This file is part of the APS files in the REVTeX 4.2 distribution.
%%   Version 4.2b of REVTeX, December 2018.
%%
%%   Copyright (c) 2019 The American Physical Society.
%%
%%   See the REVTeX 4.2 README file for restrictions and more information.
%%
\documentclass[twocolumn,secnumarabic,amssymb, nobibnotes, aps, prd]{revtex4-2}
%\usepackage{acrofont}%NOTE: Comment out this line for the release version!
\newcommand{\revtex}{REV\TeX\ }
\newcommand{\classoption}[1]{\texttt{#1}}
\newcommand{\macro}[1]{\texttt{\textbackslash#1}}
\newcommand{\m}[1]{\macro{#1}}
\newcommand{\env}[1]{\texttt{#1}}
\renewcommand{\ne}{n_{e}}  % shadow LaTeX kernel's \ne ("not equal") -- not used in this doc
\newcommand{\nfc}{n_{\mathrm{FC}}}
\newcommand{\ncx}{n_{\mathrm{CX}}}
\newcommand{\Vfc}{\lvert V_{\mathrm{FC}}\rvert}
\newcommand{\Vcx}{\lvert V_{\mathrm{CX}}\rvert}
\newcommand{\Vfcr}{V_{\mathrm{FC},r}}
\newcommand{\Vcxr}{V_{\mathrm{CX},r}}
\newcommand{\Si}{S_{i}}
\newcommand{\Scx}{S_{\mathrm{CX}}}
\newcommand{\Sreact}{S}
\newcommand{\Dped}{D_{\mathrm{ped}}}
\newcommand{\fFC}{f_{\mathrm{FC}}}
\newcommand{\fCX}{f_{\mathrm{CX}}}
\newcommand{\fsa}[1]{\left\langle #1 \right\rangle}
\newcommand{\gradrsq}{\fsa{|\nabla r|^{2}}}
\newcommand{\neginf}{x=-\infty}
\newcommand{\dneneg}{\left.\dfrac{d\ne}{dx}\right|_{\neginf}}
\newcommand{\xin}{x_{\mathrm{inner}}}
\newcommand{\gFC}{\fsa{|\nabla r|^{2}\,\nfc}}
\newcommand{\gCX}{\fsa{|\nabla r|^{2}\,\ncx}}
\newcommand{\dnexnow}{\dfrac{d\ne}{dx}}
% \newcommand{\xin}{x_{\mathrm{inner}}}
\newcommand{\nFCzero}{\fsa{\nfc}(0)}
\setlength{\textheight}{9.5in}
\usepackage{amsmath}
\usepackage{appendix}
\usepackage{hyperref}
\usepackage{graphicx}
\usepackage{bm}
\usepackage{booktabs}
\usepackage{subcaption}
% \usepackage{bm}

\begin{document}

\title{Using Firedrake to solve a reduced model for the tokamak pedestal density} 

\author{John Anthony Labbate}
\email{john.a.labbate@columbia.edu}
\affiliation{
Columbia Fusion Research Center
}
\affiliation{
Department of Applied Physics and Applied Mathematics, Columbia University, New York, New York 10027, USA
}

\date{May 2026}%

\begin{abstract}
We present a Firedrake-based finite-element implementation of an extended Saarelma–Connor tokamak pedestal electron and neutral-transport model.
The extended model is formulated as a coupled system for electron and neutral densities with flexible boundary conditions and initial guesses. 
Preliminary calculations suggest improved robustness and closer agreement with experimental density profiles than the original model implementation. 
Among the configurations examined, Dirichlet conditions on the electron density at both radial boundaries, combined with Picard iteration from experimental, piecewise-linear, or hyperbolic-tangent initial profiles yield the best match to the measured pedestal. 
The legacy solver remains substantially faster than the new implementation under current settings, with the former completing in roughly one quarter of the run time of the latter. 
Ongoing work will extend systematic comparisons between solvers and broaden validation across experimental datasets.
Additionally, timing and free parameter optimization will be explored by future work.
\end{abstract}

\maketitle


% \tableofcontents

\section{Introduction}
\label{sect:intro}
Fusion energy is the gold standard of energy production due to its immense fuel efficiency and minimal environmental impact. 
Magnetic confinement fusion is the leading energy production choice between fusion concepts. 
A magnetic confinement fusion concept known as the tokamak has shown strong promise for a power plant design due to high confinement operation.

Modern tokamak plasmas operate primarily in two modes - a lower confinement regime (L-mode) and a higher confinement regime (H-mode).
H-mode is characterized by the presence of a large transport barrier in the outer few layers of the tokamak plasma. 
This transport barrier, known as the pedestal, may only traverse a few percent of the total plasma domain \cite{snyder2011first}.
Therefore, the pedestal is characterized by large gradients in both the plasma temperature and density.

The pedestal plays a major role in the plasma performance of a fusion device. 
For example, the boundary condition imposed by this pedestal sets the temperature profile for the rest of the plasma \cite{saarelma2023testing}.
However, the pedestal is a major source of instabilities for a tokamak due to the presence of these large gradients. 
One of these instabilities is called an Edge-Localized Mode (ELM) \cite{snyder2011first}. 
ELMs are experimentally observed in most tokamaks, but at a power plant scale, ELMs are expected to be catastrophic for a fusion device \cite{liang2011overview}. 
Future fusion devices and power plants will need to operatre without ELMs to be successful.

It is clear the tokamak pedestal is of large importance. 
An understanding of the pedestal could save the tokamak from the problem of ELMs and provide better operating points for devices.
However, the pedestal is difficult to model as its large gradients pass through multiple different physics regimes.
In 2009, \cite{snyder2009development} released a model - EPED - to predict the width and height of the pressure of a tokamak pedestal.
EPED combines peeling-ballooning magnetohydrodynamic stability with Kinetic Ballooning Mode (KBM) stability to predict this pressure height and width \cite{snyder2009development}.
While EPED is a very well-studied code (for example, \cite{snyder2011first}), it has a major drawback of requiring knowledge of the density profile to predict the pressure profile.

In 2023, \cite{saarelma2023testing} released a model - the Saarelma-Connor model - that models the density profile of the tokamak pedestal to generate an input for EPED.
This is a physics-based model that will be described in detail in Sect. \ref{sect:methods}.
Unlike for EPED, a public computational resource for this Saarelma-Connor model does not yet exist.
In this study, we develop a computational resource for a modified version of the Saarelma-Connor model.
The Saarelma-Connor model makes various assumptions to generate the equation of interest for the electron density in the pedestal.
This equation is solved numerically using an implicit method.
With an implicit method, it is essential to have a good initial guess, which can be difficult to predict (although \cite{saarelma2023testing} offers a tentative solution for this).
The need for a good guess for the electron density \textit{a priori} is essential. 
However, previous work (from \cite{JohnLabbate} has illustrated that this implicit method is difficult to use to find solver convergence and accurate results that best replicate the experimental p-file.

For this reason, we take a novel approach to solving the Saarelma-Connor model.
This approach involves stripping lots of the algebra and first integration completed in \cite{saarelma2023testing} and beginning with the model's governing equations.
We implement a Piccard iteration numerical method to solve this system of equations using Firedrake and PETSc.
We compare the novel implementation of the Saarelma-Connor model with the original implicit Saarelma-Connor model through this report.


\section{A new solver for the Saarelma-Connor model}
All variable identities can be found in App. \ref{app:Variables}.

\label{sect:methods}
\subsection{The original Saarelma-Connor model}
The Saarelma-Connor model is a three-fluid model comprised of an electron fluid, a Franck-Condon (FC) monoenergetic neutral fluid, and a charge-exchange (CX) neutral fluid \cite{saarelma2023testing}.
The governing equations for these fluids in the pedestal region are given in Eq. \ref{eq:1}-\ref{eq:3}.
\begin{align}
  \nabla\cdot(\Dped\,\nabla\ne)
    &= -\,\ne\,(\nfc + \ncx)\,\Si, \label{eq:1} \\
  \nabla\cdot(V_{FC}\,\nfc)
    &= -\,\ne\,\bigl(\nfc\Si + \nfc\Scx\bigr), \label{eq:2} \\
  \nabla\cdot(V_{CX}\,\ncx)
    &= -\,\ne\,\bigl(\ncx\Si - \tfrac{1}{2}\nfc\Scx\bigr).
    \label{eq:3}
\end{align}
Eq. \ref{eq:1}-\ref{eq:3} model the diffusive steady-state motion of each fluid with source and sink terms. 
In Eq. \ref{eq:1}, electron impact ionization of both FC and CX neutrals provide an electron source.
There are no electron sinks in this model.
In Eq. \ref{eq:2}, electron impact ionization of FC neutrals and charge exchange with a bulk-plasma ion (assuming plasma quasi-neutrality, so the electron and ion densities are locally equal) are both FC neutral sinks.
In Eq. \ref{eq:3}, electron impact ionization of CX neutrals provide a CX neutral sink, and charge exchange with FC neutrals provides a CX neutral source. 
The CX neutral source term has a $\frac{1}{2}$ coefficient because it assumes these particles have random trajectories, so half of these neutrals will be ballistically lost to the wall of the device.

In the Saarelma-Connor model, many assumptions regarding the pedestal region are made.
Ultimately, the system is reduced to a single 1D ordinary differential equation for the electron density using flux surface averaging, which is defined as $\langle A \rangle = \oint R^2Ad\theta / \oint R^2 d\theta$ for some quantity $A$.
To obtain this 1D equation, an analytic integration is completed.
The full derivation of the final equation in the Saarelma-Connor model can be found in \cite{saarelma2023testing} and \cite{saarelma2024density}, but the final equations are:
\begin{align}
\frac{d}{dx}\left(L_2(n_e)\right)
=
\frac{n_e S_i}{|V_{FC}|}
\left(1+\frac{S_{CX}}{S_i}\right)
L_2(n_e),
% \tag{12}
\end{align}

\begin{widetext}
\begin{align}
    \left[
1-
\frac{|V_{FC}|f_{FC}}{|V_{CX}|f_{CX}}
\left(
\frac{S_i+\frac{S_{CX}}{2}}{S_i+S_{CX}}
\right)
\right]
L_2(n_e)
=
-\frac{C}{|V_{CX}|f_{CX}}
+
\frac{1}{|V_{CX}|f_{CX}}
D_{\mathrm{ped}}
\frac{dn_e}{dx}
-
\frac{1}{n_e S_i}
\frac{d}{dx}
\left(
\langle |\nabla r|^2 \rangle
D_{\mathrm{ped}}
\frac{dn_e}{dx}
\right).
% \tag{13}
% \]
\end{align}
\end{widetext}
\cite{saarelma2023testing} suggests utilizing an implicit solving method due to the stiffness of this third-order equation. 
The implicit solver works by first solving a simplified model for $n_e$ where $n_{CX}=0$ is assumed (Eq. \ref{eq:16}).
Then, the full model is solved using Eq. \ref{eq:15-derived}, but the $n_e$ within the integral is assumed to be the $n_e$ found in Eq. \ref{eq:16}.
\begin{widetext}

\begin{equation}
    \frac{d}{dx}\!\left[\gradrsq\,\Dped\,\frac{d\ne}{dx}\right]
    \;=\; \ne\,\Dped\,\frac{\Si+\Scx}{\Vfc}\,
          \left(\frac{d\ne}{dx} - \dneneg\right).
  \label{eq:16}
\end{equation}

\begin{equation}
  \begin{aligned}
    \frac{d}{dx}\!\left[\gradrsq\,\Dped\,\dnexnow\right]
      &= -\,\ne\,\Si\,\Biggl\{
            -\,\frac{\gradrsq\,\Dped}{\Vcx\,\fCX}\bigl(\dnexnow-\dneneg\bigr)  \\
      &\qquad + \left[\,1 - \frac{\Vfc\,\fFC(\Si+\Scx/2)}{\Vcx\,\fCX(\Si+\Scx)}\,\right]
                \nFCzero\,\exp\!\left[\int_{0}^{x}\!\frac{\ne(x')(\Si+\Scx)}{\fFC\,\Vfc}\,dx'\right]
         \Biggr\}.
  \end{aligned}
  \label{eq:15-derived}
\end{equation}
\end{widetext}
Eq. \ref{eq:15-derived} is the form we derived, which includes factors that are missing from \cite{saarelma2023testing}.

While this solution method may seem advantageous, it relies heavily on the knowledge of the large free parameter space that is included in this model \cite{JohnLabbate}.
Specifically, this implicit method has a difficult time converging when implemented with Scipy's \texttt{solve\_bvp}.
This solution method also introduces additional assumptions to the original model, which already is not fully capturing the behavior of the tokamak pedestal (diffusion equations are really a proxy for more complex transport and turbulence occurring on smaller scales).

\subsection{A modified Saarelma-Connor model}
In this study, we relax the assumptions imposed by \cite{saarelma2023testing} and look for a more cooperative numerical solver by solving a flux-averaged version of the inital coupled diffusion equations.
By simply taking the flux surface average of Eq. \ref{eq:1}-\ref{eq:3}, the following are obtained.
\begin{equation}
  \frac{d}{dx}\!\left[\gradrsq\,\Dped\,\frac{d\ne}{dx}\right]
    \;=\; -\,\ne\,\Si\,\bigl(\fsa{\nfc}+\fsa{\ncx}\bigr),
  \label{eq:8}
\end{equation}
\begin{align}
  \Vfc\,\frac{d}{dx}\!\bigl[\fFC\,\fsa{\nfc}\,&\gradrsq\bigr]
    \;=\; 
    \notag \\ &\ne\,(\Si+\Scx)\,\fsa{\nfc},
      \label{eq:9}
\end{align}
\begin{align}
  \Vcx\,\frac{d}{dx}\!\bigl[\fCX\,&\fsa{\ncx}\,\gradrsq\bigr]
    \;=\; 
    \notag \\ &\ne\left(\fsa{\ncx}\Si - \frac{\Scx}{2}\fsa\nfc\right).
      \label{eq:10}
\end{align}
This is a coupled system of ordinary differential equations (ODEs), one second-order and two first-order.
We now can directly solve this system using Firedrake with no additional assumptions imposed.

\subsection{Firedrake implementation of coupled system}
The implementation of Eq. \ref{eq:16} and \ref{eq:15-derived} can be found in \url{https://github.com/jlabbate15/saarelma-conner-ped} in the solver.py file. 
The Firedrake implementation of Eq. \ref{eq:8}-\ref{eq:10} exists in the same repository but in the solver-firedrake.py file.
Both methods use the same import structure, which will not be described in this report as it is very tedious and incorporates only fixed engineering parameters.
Interestingly, the Eq. \ref{eq:16} and \ref{eq:15-derived} is much more robust when utilizing non-dimensionalization, but this is not required of the Firedrake implementation of Eq. \ref{eq:8}-\ref{eq:10}.

To solve a system of coupled ordinary differential equations using Firedrake, we do the following. 
First, we must determine which boundary conditions to use.
Each 1D ODE requires a single boundary condition, and the 2D ODE requires two boundary conditions.
For $n_{CX}$ and $n_{FC}$, the obvious choice is the Dirichlet boundary condition at the outer-most part of the pedestal (the separatrix), as prior knowledge of these values is much more likely than the inner boundary condition. 
For $n_e$, the Dirichlet outer boundary condition is also a safe choice. 
For the inner boundary condition of $n_e$, however, a choice can be made between a Dirichlet and a Neumann boundary condition. 
Choosing a Neumann boundary condition corresponds slightly better to the physics occuring at this point in the pedestal, which is dominated by electron flux (proportional to the first derivative of the electron density).
However, the Neumann boundary condition is not as strict of a constraint on the pedestal profile as the Dirichlet boundary condition.
We will compare both options in this study.

Because this is a coupled system, the terms we are solving for appear in the differential equations themselves. 
Therefore, we need to use an iterative method to solve this. 
We will use a Picard (or fixed-point) iteration method to tackle this system. 
When using Picard iteration, the unknowns that are not differentiated are "lagged" behind the others, meaning their value from the previous step is used to solve for the current step's value via the differential equation of interest. 
Therefore, a starting point for these unknowns is required.
At first, providing an initial guess may seem problematic if we need knowledge of the densities prior to the solve. 
However, we will show that for all tested guesses, even those with no experimental or physical knowledge, the solver will converge on a single solution no matter the initial guess. 
Current analysis suggests the original Saarelma-Connor implementation does not converge to equivalent solutions when varying the initial guess \cite{JohnLabbate}.
The following initial guesses will be tested.
\begin{itemize}
    \item Experimental profile.
    \item Linear profile between boundary conditions.
    \item Hyperbolic tangent profile.
\end{itemize}
The hyperbolic tangent profile is commonly used to model the tokamak edge, which is why it serves as a good initial guess for this model.

This is a non-linear system, since the three unknowns are found to multiply each other in multiple terms.
Therefore, the essential step in solving this system is the development of the residual. 
The residual is developed by writing all three coupled equations in form $F_i(x)=0$, where $i$ is the equation number.
Then, each residual is multiplied by a test function and integrated once over the $x$ domain to obtain the weak form of each equation to form the weak form of these equations. 
The weak form has solutions that will correspond to the strong form of these differential equations (the strong form is the non-test function and integrated version).
The development of the weak forms of these equations can be found in App. \ref{app:WeakForms}. 
The resulting residuals are summed and then fed into the Picard iteration.

A solver is specified for the Picard iterations. 
SNES implements a Newton method with linesearch for the non-linear solver.
Since these are one-dimensional equations, an LU factorization is computationally suitable for all linear solves. 
An element of future work is to implement multigrid schemes.
Further combinations of solvers are possible and may yield faster results, however, the current configuration is suitable for this study.


\section{Comparison of solution methods}
We successfully implemented the modified coupled Saarelma-Connor model for the tokamak pedestal in this study.
Unfortunately, there is no good way to compare these outputs and verify there are no bugs in the code. 
The best we can do is to compare to the experimental outputs, knowing that there will be a difference as this Saarelma-Connor model is not capturing all of the physics.
In fact, \cite{saarelma2023testing} shows significant spreads of the pedestal information in experiments and predicted by their model, so we should expect similar results.

To check the preliminary success of this model over its various setup options, a scan of the different available boundary conditions at the inner boundary for $n_e$ and the different initial guesses was completed.
The results of this scan are shown in Fig. \ref{fig:1}.
\begin{figure}[htbp]
  \centering
  \begin{subfigure}[b]{0.45\textwidth}
    \centering
    \includegraphics[width=\textwidth]{ne_scan.png}
    \caption{$n_e$}
    \label{fig:sub1}
  \end{subfigure}
  \hfill
  \begin{subfigure}[b]{0.45\textwidth}
    \centering
    \includegraphics[width=\textwidth]{nfc_scan.png}
    \caption{$n_{FC}$}
    \label{fig:sub2}
  \end{subfigure}
    \begin{subfigure}[b]{0.45\textwidth}
    \centering
    \includegraphics[width=\textwidth]{ncx_scan.png}
    \caption{$n_{CX}$}
    \label{fig:sub3}
  \end{subfigure}
    \caption{Pedestal profiles of the Saarelma-Connor model species for DIII-D shot 15840 at time 03000 milliseconds. Illustrating scan over boundary condition type at inner boundary for $n_e$ and different initial guesses. Boundary conditions are always calculated from the p-file (experimental data). Note p-files (experimental data) do not contain information about $n_{CX}$ and $n_{FC}$, so this result only appears on the $n_e$ plot.}
    \label{fig:1}
\end{figure}
It is important to note that the modified Firedrake-implemented Saarelma-Connor model converges regularly and has not been subject to optimization of its free parameters.
The original Saarelma-Connor model struggles to converge (a converged case is depicted in Fig. \ref{fig:1}) and scans a 5D free parameter space to find an optimal parameter combination based on how far the result of the model is from the original p-file (experimental) input.
This optimal combination is plotted in Fig. \ref{fig:1} as "Original Saarelma-Connor model".
It is surprising to find that even with a lack of optimization, the Firedrake-implemented model outperforms the original model quite significantly.

Fig. \ref{fig:1} also indicates that for constant boundary condition, any of the initial guesses given to the Firedrake model implementation result in the same model answer. 
It is the boundary condition at the inner boundary that shifts the solution, and currently this case finds the Dirichlet boundary condition to match the experiment much better.
With a Dirichlet boundary condition at the inner pedestal boundary, it is likely not guaranteed that the Neumann condition would match the experiments because this is the pre-pedestal region where the physics are not controlled by the same assumptions that are present in the Saarelma-Connor model.

Future work should optimize the free parameters, including the placement of the inner boundary condition, of the Firedrake-implemented Saarelma-Connor model.

Future work should also continue to search for bugs in this code's implementation and explore its effectiveness by applying this model to many tokamak pedestal experimental cases and producing a plot similar to the validation work completed in \cite{saarelma2023testing}.

Short timing tests were also completed for the Firedrake implementation of the Saarelma-Connor model. 
These tests were completed on the scans that generated Fig. \ref{fig:1}.
The original model implementation in Scipy's \texttt{solve\_bvp} runs in about 0.08 seconds, and the Firedrake implementation runs in about 0.30 seconds.
The original model has already been optimized for time, however, where the Firedrake model has not. 
Either way, both times are sufficient for larger scans of both model implementations.


\section{Conclusions}
In this study, we implemented a modified, more general version of the Saarelma-Connor model \cite{saarelma2023testing} using Firedrake to build a custom numerical solver.
Preliminary results indicate more robust convergence and stronger matches to experiment for the Firedrake implementation over the original Saarelma-Connor model.
A few different Firedrake configurations were tested for the new solver.
It was found that Dirichlet boundary conditions everywhere in the system and any initial guess for a Picard iteration (either from the experiment, linear fit between the experimental boundary points, or a hyperbolic tangent profile) provided the closest match to the experimental result.
The timing-optimized original model implementation was found to be about four times faster than the non-timing-optimized novel implementation.



\begin{acknowledgments}
Thank you for great semester!
\end{acknowledgments}


\bibliographystyle{apsrev4-2} % Sets the PRL-specific formatting
\bibliography{bib} % Points to my_plasma_refs.bib


\appendix
\section{Weak form derivation for coupled system}
\label{app:WeakForms}
This appendix documents the variational formulation implemented in
\texttt{solver-firedrake.py} for the modified coupled Saarelma--Connor pedestal model~\cite{saarelma2023testing}.
The three unknowns are flux-surface-averaged profiles on the radial
coordinate $x = r - r_{\mathrm{sep}}$:
\begin{equation}
  n_e(x),\quad \langle n_{\mathrm{FC}}\rangle(x),\quad \langle n_{\mathrm{CX}}\rangle(x),
\end{equation}
with $x=0$ at the separatrix and $x<0$ in the pedestal.
Coefficients $D_{\mathrm{ped}}(x)$, $S_i(x)$, $S_{\mathrm{CX}}(x)$,
$\langle|\nabla r|^{2}\rangle(x)$, and $|V_{\mathrm{CX}}|(x)$ are precomputed
on the parent grid and interpolated onto the Firedrake mesh;
$|V_{\mathrm{FC}}|$ and the form factors $f_{\mathrm{FC}}$, $f_{\mathrm{CX}}$
are treated as constants in the residual.

\subsection{Domain, function spaces, and boundary conditions}
\label{sec:domain}

The computational domain is the closed interval
\begin{equation}
  \Omega = [x_{\mathrm{inner}},\,0], \qquad x_{\mathrm{inner}} < 0.
\end{equation}
A one-dimensional \texttt{IntervalMesh} is used with boundary labels
\texttt{ds(1)} at $x=x_{\mathrm{inner}}$ (inner pedestal foot) and
\texttt{ds(2)} at $x=0$ (separatrix; inflow for neutrals).

The mixed finite-element space is
\begin{equation}
  W = V \times V \times V, \qquad
  u = \bigl(n_e,\,\langle n_{\mathrm{FC}}\rangle,\,\langle n_{\mathrm{CX}}\rangle\bigr) \in W,
\end{equation}
with $V = \mathrm{CG}_p(\Omega)$ (default $p=2$) where $CG_p$ is a continuous Galerkin space with order $p$.
Test functions are denoted $v_e, v_F, v_C \in V$.

Four scalar boundary conditions close the fourth-order coupled system:
\begin{subequations}
\label{eq:bcs}
\begin{align}
  n_e(0) &= n_{e,0}, \label{eq:bc-ne0}\\
  \langle n_{\mathrm{FC}}\rangle(0) &= \langle n_{\mathrm{FC}}\rangle_{0}, \label{eq:bc-nfc0}\\
  \langle n_{\mathrm{CX}}\rangle(0) &= \langle n_{\mathrm{CX}}\rangle_{0}, \label{eq:bc-ncx0}
\end{align}
\end{subequations}
with Dirichlet data at the separatrix, and at $x=x_{\mathrm{inner}}$ either
\begin{equation}
  n_e(x_{\mathrm{inner}}) = n_{e,\mathrm{inner}}
  \quad\text{(Dirichlet)}
  \label{eq:bc-ne-dirichlet}
\end{equation}
or
\begin{equation}
  \left.\frac{dn_e}{dx}\right|_{x=x_{\mathrm{inner}}} = n'_{e,\mathrm{inner}}
  \quad\text{(Neumann / flux)}.
  \label{eq:bc-ne-neumann}
\end{equation}
Essential Dirichlet conditions are enforced with \texttt{DirichletBC}.
% The Neumann condition for $n_e$ enters as a natural boundary term
% (App. A.\ref{sec:ne-weak}).

All strong forms are given by Eq. \ref{eq:8}-\ref{eq:10}. \\

% \subsection{Strong forms solved}
% \label{sec:strong}

% After flux-surface averaging and introduction of the form factors
% (Eq.~(7) of Ref.~\cite{saarelma2023testing}), the code solves the following system on
% $\Omega$:
% \begin{subequations}
% \label{eq:strong-system}
% \begin{align}
%   \frac{d}{dx}\!\left[\langle|\nabla r|^{2}\rangle\,D_{\mathrm{ped}}\,\frac{dn_e}{dx}\right]
%     &= -\,n_e\,S_i\,(\langle n_{\mathrm{FC}}\rangle + \langle n_{\mathrm{CX}}\rangle),
%   \label{eq:strong-ne}\\[4pt]
%   |V_{\mathrm{FC}}|\,\frac{d}{dx}\!\left[f_{\mathrm{FC}}\,\langle n_{\mathrm{FC}}\rangle\,\langle|\nabla r|^{2}\rangle\right]
%     &= +\,n_e\,(S_i + S_{\mathrm{CX}})\,\langle n_{\mathrm{FC}}\rangle,
%   \label{eq:strong-nfc}\\[4pt]
%   |V_{\mathrm{CX}}|\,\frac{d}{dx}\!\left[f_{\mathrm{CX}}\,\langle n_{\mathrm{CX}}\rangle\,\langle|\nabla r|^{2}\rangle\right]
%     &= n_e\,\left(\langle n_{\mathrm{CX}}\rangle\,S_i - \tfrac{1}{2}\,S_{\mathrm{CX}}\,\langle n_{\mathrm{FC}}\rangle\right).
%   \label{eq:strong-ncx}
% \end{align}
% \end{subequations}
% Equations~\eqref{eq:strong-nfc} and~\eqref{eq:strong-ncx} are written in
% conservative form for the fluxes
% $F_F \equiv |V_{\mathrm{FC}}|\,f_{\mathrm{FC}}\,\langle n_{\mathrm{FC}}\rangle\,\langle|\nabla r|^{2}\rangle$
% and
% $F_C \equiv |V_{\mathrm{CX}}|\,f_{\mathrm{CX}}\,\langle n_{\mathrm{CX}}\rangle\,\langle|\nabla r|^{2}\rangle$.
% All quantities are in physical (SI) units; no non-dimensionalization
% is applied in the Firedrake residual.

\subsection{Electron density weak form}
\label{sec:ne-weak}

% \paragraph{Weighted residual.}
Multiply Eq.~\eqref{eq:8} by a test function $v_e \in V$ and
integrate over $\Omega$:
\begin{align}
  \int_{x_{\mathrm{inner}}}^{0}
    \frac{d}{dx}\!&\left[\langle|\nabla r|^{2}\rangle\,D_{\mathrm{ped}}\,\frac{dn_e}{dx}\right] v_e \, dx
  \notag \\ &= -\int_{x_{\mathrm{inner}}}^{0} n_e\,S_i\,(\langle n_{\mathrm{FC}}\rangle + \langle n_{\mathrm{CX}}\rangle)\, v_e \, dx.
  \label{eq:ne-weighted}
\end{align}

% \paragraph{Integration by parts.}
% Define the diffusive flux
% $q \equiv \langle|\nabla r|^{2}\rangle\,D_{\mathrm{ped}}\,n_e'$.
Now integrate by parts.
With primes denoting $d/dx$,
\begin{widetext}
\begin{align}
  -\int_{x_{\mathrm{inner}}}^{0} n_e\,S_i\,(\langle n_{\mathrm{FC}}\rangle + \langle n_{\mathrm{CX}}\rangle)\, v_e \, dx =   \underbrace{\Bigl[\langle|\nabla r|^{2}\rangle\,D_{\mathrm{ped}}\,n_e'\, v_e\Bigr]_{x=x_{\mathrm{inner}}}^{x=0}}_{\text{boundary terms}}
  - \int_{x_{\mathrm{inner}}}^{0} \langle|\nabla r|^{2}\rangle\,D_{\mathrm{ped}}\,n_e'\, v_e' \, dx
  \label{eq:ne-ibp}
\end{align}
\end{widetext}

% \paragraph{Boundary terms.}
At any boundary where $n_e$ is prescribed (Dirichlet), the test function satisfies $v_e=0$.
At $x=0$, $n_e(0)=n_{e,0}$ is Dirichlet, so the boundary contribution at $x=0$ vanishes.
At $x=x_{\mathrm{inner}}$ with Dirichlet data~\eqref{eq:bc-ne-dirichlet}, $v_e(x_{\mathrm{inner}})=0$ and the inner boundary term also vanishes.
At $x=x_{\mathrm{inner}}$ with Neumann data~\eqref{eq:bc-ne-neumann},
$v_e(x_{\mathrm{inner}})$ is not constrained and the prescribed gradient
enters through
\begin{equation}
  % \Bigl[\langle|\nabla r|^{2}\rangle\,D_{\mathrm{ped}}\,n_e'\, v_e\Bigr]_{x=x_{\mathrm{inner}}}
  % \longrightarrow
  \langle|\nabla r|^{2}\rangle(x_{\mathrm{inner}})\,D_{\mathrm{ped}}(x_{\mathrm{inner}})\,n'_{e,\mathrm{inner}}\, v_e(x_{\mathrm{inner}}).
  \label{eq:ne-neumann-natural}
\end{equation}

% \paragraph{Volume weak form.}
Collecting terms, the semi-discrete residual for Eq.~\eqref{eq:8}
(without Picard lagging) is: find $n_e \in V$ such that, for all admissible
$v_e \in V$ where $\delta_{\mathrm{N}}=1$ for a Neumann inner BC and
$\delta_{\mathrm{N}}=0$ for Dirichlet, 
\begin{widetext}
\begin{equation}
  \int_{x_{\mathrm{in}}}^{0} \langle|\nabla r|^{2}\rangle\,D_{\mathrm{ped}}\,n_e'\, v_e' \, dx
  - \int_{x_{\mathrm{in}}}^{0} n_e\,S_i\,(\langle n_{\mathrm{FC}}\rangle + \langle n_{\mathrm{CX}}\rangle)\, v_e \, dx
  + \delta_{\mathrm{N}}\,
    \langle|\nabla r|^{2}\rangle(x_{\mathrm{in}})\,D_{\mathrm{ped}}(x_{\mathrm{in}})\,n'_{e,\mathrm{in}}\, v_e(x_{\mathrm{in}})
  = 0.
  \label{eq:weak-ne}
\end{equation}
\end{widetext}

% In Firedrake, the volume terms are
% \texttt{(g\_fd * D\_fd * ne.dx(0) * v\_e.dx(0) - ne * Si\_fd * (nfc + nCX) * v\_e) * dx}
% and the Neumann contribution is
% \texttt{g\_fd * D\_fd * dne\_dx\_inner * v\_e * ds(1)}.

\subsection{Franck-Condon neutrals residual}
\label{sec:nfc-weak}

% Equation~\eqref{eq:9} is first order with characteristics
% directed from the separatrix inward.
% The implementation uses a strong-form Galerkin residual (no integration by parts on the flux derivative), so that the discrete solution matches a forward march from the inflow boundary $x=0$.
Equation~\eqref{eq:9} is first order and does not require integration like parts, unlike the electron neutral weak form derivation.
Integration by parts would transfer the derivative onto $v_F$ and introduce an inflow boundary term that does not preserve the upwind structure.
Multiplying by $v_F \in V$ and integrating gives the following, for all $v_F \in V$:
% and with $|V_{\mathrm{FC}}|$ and $f_{\mathrm{FC}}$ are spatially constant in the code.
\begin{widetext}
\begin{align}
  \int_{x_{\mathrm{inner}}}^{0}
    \left\{
      |V_{\mathrm{FC}}|\,\frac{d}{dx}\!\left[f_{\mathrm{FC}}\,\langle n_{\mathrm{FC}}\rangle\,\langle|\nabla r|^{2}\rangle\right]
      - n_e\,(S_i + S_{\mathrm{CX}})\,\langle n_{\mathrm{FC}}\rangle
    \right\} v_F \, dx = 0,
  \label{eq:weak-nfc}
\end{align}
% \begin{equation}
%   |V_{\mathrm{FC}}|\,\frac{d}{dx}\!\left[f_{\mathrm{FC}}\,\langle n_{\mathrm{FC}}\rangle\,\langle|\nabla r|^{2}\rangle\right]
%   = \frac{d}{dx}\!\left[|V_{\mathrm{FC}}|\,f_{\mathrm{FC}}\,\langle|\nabla r|^{2}\rangle\,\langle n_{\mathrm{FC}}\rangle\right],
% \end{equation}
\end{widetext}

% which corresponds to \texttt{(VFC\_const * fFC\_const * g\_fd * nFC).dx(0) * v\_F * dx}.

\subsection{Charge-exchange neutrals residual}
\label{sec:ncx-weak}

The same reasoning as the Franck-Condon neutrals applies to Eq.~\eqref{eq:10}.
Thus, for all $v_C \in V$.
\begin{widetext}
\begin{align}
  \int_{x_{\mathrm{inner}}}^{0}
    \left\{
      |V_{\mathrm{CX}}|\,\frac{d}{dx}\!\left[f_{\mathrm{CX}}\,\langle n_{\mathrm{CX}}\rangle\,\langle|\nabla r|^{2}\rangle\right]
      - n_e\,\left(\langle n_{\mathrm{CX}}\rangle\,S_i - \tfrac{1}{2}\,S_{\mathrm{CX}}\,\langle n_{\mathrm{FC}}\rangle\right)
    \right\} v_C \, dx = 0.
  \label{eq:weak-ncx}
\end{align}
\end{widetext}

\subsection{Coupled residual and Picard linearization}
\label{sec:coupled}

The total residual solved by Firedrake is
\begin{equation}
  F(u;\,v_e,v_F,v_C) = F_1 + F_2 + F_3 = 0,
  \label{eq:total-F}
\end{equation}
where $F_1$, $F_2$, $F_3$ are the forms derived above, plus the optional
Neumann boundary term in $F_1$.
Each nonlinear term is a product of two unknowns
(e.g.\ $n_e\,S_i\,\langle n_{\mathrm{FC}}\rangle$).
With \texttt{use\_picard=True}, one factor is lagged at iterate $u^{(k)}$
while solving for $u^{(k+1)}$:
$F_1$ lags $(\langle n_{\mathrm{FC}}\rangle+\langle n_{\mathrm{CX}}\rangle)$,
$F_2$ lags $n_e$, and $F_3$ lags $n_e$ and $\langle n_{\mathrm{FC}}\rangle$
in the CX source.
With \texttt{use\_picard=False}, all products use the current iterate and a single SNES Newton solve is applied to the fully coupled residual (this option is not explored in this study).
Updates may be under-relaxed with the control of the user,
$u^{(k+1)} \leftarrow (1-\lambda)\,u^{(k)} + \lambda\,\tilde u^{(k+1)}$,
with $\lambda=\texttt{relax}$ (default $0.5$).

% \subsection{Mapping to \texorpdfstring{\texttt{solver-firedrake.py}}{solver-firedrake.py}}
% \label{sec:code-map}

\begin{table*}[t]
\caption{\label{tab:code-map}Correspondence between strong forms, variational residuals, and Firedrake symbols.}
\begin{ruledtabular}
\begin{tabular}{lll}
Strong form & Weak / Galerkin residual & Code symbol \\
\hline
Eq.~\eqref{eq:8} & Eq.~\eqref{eq:weak-ne} & \texttt{F1} (+ \texttt{F1\_neumann\_bc}) \\
Eq.~\eqref{eq:9} & Eq.~\eqref{eq:weak-nfc} & \texttt{F2} \\
Eq.~\eqref{eq:10} & Eq.~\eqref{eq:weak-ncx} & \texttt{F3} \\
\end{tabular}
\end{ruledtabular}
\end{table*}

% Essential boundary conditions~\eqref{eq:bc-ne0}--\eqref{eq:bc-ncx0} and
% optional Dirichlet data~\eqref{eq:bc-ne-dirichlet} are enforced via
% \texttt{DirichletBC}.
% The Neumann condition~\eqref{eq:ne-neumann-natural} is the only
% $n_e$-related boundary term appearing explicitly in the residual.


\subsection{Rederiving electron density weak form}
\label{sec:ne-weak-gen}

We extend the baseline formulation to incorporate a highly non-linear, physics-based transport model where the total particle diffusion coefficient $D$ combines a turbulent pedestal contribution $\Dped$ and a neoclassical baseline contribution $D_{\mathrm{NEO}}(x)$:
\begin{equation}
D = \Dped + D_{\mathrm{NEO}}(x)
\end{equation}
The turbulent pedestal transport incorporates Electron Temperature Gradient (ETG) scale anomalies and Kinetic Ballooning Mode (KBM) scale instabilities, such that $\Dped = D_{\mathrm{ETG}} + D_{\mathrm{KBM}}$. Following their known scalings, we define:
\begin{equation}
D_{\mathrm{ETG}} = \frac{C_{\mathrm{ETG}}(x)}{\ne}, \qquad D_{\mathrm{KBM}} = C_{\mathrm{KBM}}(x) \frac{d}{dx}(\ne T_e)
\end{equation}
where $C_{\mathrm{ETG}}(x)$ and $C_{\mathrm{KBM}}(x)$ represent spatial calibration profile parameters and $T_e(x)$ is the precomputed, known electron temperature background profile. Applying the product rule to the KBM gradient term gives:
\begin{equation}
D_{\mathrm{KBM}} = C_{\mathrm{KBM}}(x) \left( T_e \frac{d\ne}{dx} + \ne \frac{dT_e}{dx} \right)
\end{equation}
This yields a fully state-dependent definition for the complete transport diffusion coefficient:
\begin{widetext}
\begin{equation}
D\left(\ne, \frac{d\ne}{dx}\right) = \frac{C_{\mathrm{ETG}}(x)}{\ne} + C_{\mathrm{KBM}}(x) T_e \frac{d\ne}{dx} + C_{\mathrm{KBM}}(x) \ne \frac{dT_e}{dx} + D_{\mathrm{NEO}}(x)
\end{equation}
\end{widetext}

Substituting this total diffusion framework directly into the flux-surface-averaged electron transport equation (Eq. \ref{eq:8}) produces the expanded strong form equation:
\begin{widetext}
\begin{equation}
\frac{d}{dx}\left[\gradrsq \left( \frac{C_{\mathrm{ETG}}(x)}{\ne} + C_{\mathrm{KBM}}(x) T_e \frac{d\ne}{dx} + C_{\mathrm{KBM}}(x) \ne \frac{dT_e}{dx} + D_{\mathrm{NEO}}(x) \right) \frac{d\ne}{dx}\right] = -\ne\Si\bigl(\fsa{\nfc}+\fsa{\ncx}\bigr)
\end{equation}
\end{widetext}

To form the weak statement, we multiply the strong form by an arbitrary test function $v_e \in V$ and integrate across the spatial interval domain $\Omega = [\xin, 0]$:
\begin{widetext}
\begin{equation}
\int_{\xin}^{0} \frac{d}{dx}\left[\gradrsq D\left(\ne \frac{d\ne}{dx}\right) \frac{d\ne}{dx}\right] v_e  dx = -\int_{\xin}^{0} \ne\Si\bigl(\fsa{\nfc}+\fsa{\ncx}\bigr) v_e dx
\end{equation}
\end{widetext}

We perform integration by parts on the left-hand side to shift the order of differentiation onto the test function space:
\begin{widetext}
\begin{equation}
\left.\left[\gradrsq D\left(\ne \frac{d\ne}{dx}\right) \frac{d\ne}{dx} v_e\right]\right|_{\xin}^{0} - \int_{\xin}^{0} \gradrsq D\left(\ne, \frac{d\ne}{dx}\right) \frac{d\ne}{dx} v_e'  dx = -\int_{\xin}^{0} \ne\Si\bigl(\fsa{\nfc}+\fsa{\ncx}\bigr) v_e  dx
\end{equation}
\end{widetext}

We invoke boundary constraints as specified in Sec. \ref{sec:domain}:
\begin{itemize}
\item At the separatrix region boundary ($x=0$), a fixed Dirichlet constraint $\ne(0) = n_{e,0}$ dictates that test functions vanish here ($v_e(0) = 0$), turning the upper boundary term to zero.
\item At the inner domain footprint boundary ($x=\xin$), the behavior is dictated by the configuration flag $\delta_{\mathrm{N}}$ ($\delta_{\mathrm{N}}=1$ for a Neumann scenario, $\delta_{\mathrm{N}}=0$ for a Dirichlet scenario). Under a Neumann setting, $\left.\frac{d\ne}{dx}\right|_{\xin} = n'_{e,\mathrm{inner}}$ is explicitly substituted into the residual.
\end{itemize}

Multiplying the expression by $-1$, bringing terms onto the left-hand side, and fully expanding the non-linear transport coefficient terms within the volume integral gives the final weak form residual equation. Find $\ne \in V$ such that for all $v_e \in V$:
\begin{widetext}
\begin{align}
& \int_{\xin}^{0} \gradrsq \left[ \frac{C_{\mathrm{ETG}}(x)}{\ne} \frac{d\ne}{dx} + C_{\mathrm{KBM}}(x) T_e \left(\frac{d\ne}{dx}\right)^2 + C_{\mathrm{KBM}}(x) \ne \frac{dT_e}{dx} \frac{d\ne}{dx} + D_{\mathrm{NEO}}(x) \frac{d\ne}{dx} \right] v_e'  dx 
\notag \\ & - \int_{\xin}^{0} \ne\Si\bigl(\fsa{\nfc}+\fsa{\ncx}\bigr) v_e dx + \delta_{\mathrm{N}}\gradrsq(\xin)\left[\Dped(\xin) + D_{\mathrm{NEO}}(\xin)\right] n'_{e,\mathrm{inner}} v_e(\xin) = 0
\end{align}
\end{widetext}


\subsection{Rederiving electron density weak form with total pressure}
\label{sec:ne-weak-total-pressure}

We extend the formulation to incorporate a total pressure gradient-driven transport model. The total particle diffusion coefficient $D$ combines a turbulent pedestal contribution $\Dped$ and a neoclassical baseline contribution $D_{\mathrm{NEO}}(x)$:
\begin{equation}
D = \Dped + D_{\mathrm{NEO}}(x) = D_{KBM} + D_{ETG} + D_{NEO}
\end{equation}
The turbulent pedestal transport incorporates Electron Temperature Gradient (ETG) scale anomalies and Kinetic Ballooning Mode (KBM) scale instabilities. The KBM term is driven by the total pressure gradient, $dp/dx$. Assuming quasineutrality such that the ion density $n_i \approx \ne$, the total pressure is $p = \ne T_e + n_i T_i \approx \ne (T_e + T_i)$. Following their known scalings, we define:
\begin{equation}
D_{\mathrm{ETG}} = \frac{C_{\mathrm{ETG}}(x)}{\ne}, \qquad D_{\mathrm{KBM}} = A_{KBM}+B_{\mathrm{KBM}} \frac{dp}{dx}
\end{equation}
where
\begin{align}
    A_{KBM} = -\alpha_{crit}G_{KBM},
\end{align}
\begin{align}
    B_{KBM} = \alpha_{no~p'} G_{KBM},
\end{align}
\begin{align}
    G_{KBM} = C_{KBM} \frac{\rho^2 c_s}{a},
\end{align}
\begin{align}
    \alpha_{no~p'} = -\frac{2 \partial_\psi V}{(2\pi)^2}\mu_0\sqrt{\frac{V}{2\pi^2R_0}}.
\end{align}
Applying the product rule to the total pressure gradient yields:
\begin{align}
&D_{\mathrm{KBM}} = A_{KBM}+B_{\mathrm{KBM}}(x) \Big[ (T_e + T_i) \frac{d\ne}{dx} 
\notag \\ & + \ne \frac{d}{dx}(T_e + T_i) \Big]
\end{align}
where $T_e(x)$ and $T_i(x)$ are precomputed, known background temperature profiles. The total state-dependent diffusion coefficient becomes:
\begin{widetext}
\begin{equation}
D\left(\ne, \frac{d\ne}{dx}\right) = \frac{C_{\mathrm{ETG}}(x)}{\ne} + A_{KBM} + B_{\mathrm{KBM}}(x) (T_e + T_i) \frac{d\ne}{dx} + B_{\mathrm{KBM}}(x) \ne \frac{d(T_e + T_i)}{dx} + D_{\mathrm{NEO}}(x)
\end{equation}
\end{widetext}

Substituting this total diffusion framework directly into the flux-surface-averaged electron transport equation (Eq. \ref{eq:8}) produces the expanded strong form equation:
\begin{widetext}
\begin{equation}
\frac{d}{dx}\left[\gradrsq D\left(\ne, \frac{d\ne}{dx}\right) \frac{d\ne}{dx}\right] = -\ne\Si\bigl(\fsa{\nfc} + \fsa{\ncx}\bigr)
\end{equation}
\end{widetext}

To form the weak statement, we multiply the strong form by an arbitrary test function $v_e \in V$ and integrate across the spatial interval domain $\Omega = [\xin, 0]$:
\begin{widetext}
\begin{equation}
\int_{\xin}^{0} \frac{d}{dx}\left[\gradrsq D\left(\ne, \frac{d\ne}{dx}\right) \frac{d\ne}{dx}\right] v_e  dx = -\int_{\xin}^{0} \ne\Si\bigl(\fsa{\nfc} + \fsa{\ncx}\bigr) v_e dx
\end{equation}
\end{widetext}

We perform integration by parts on the left-hand side to shift the spatial derivative onto the test function space, utilizing the standard evaluation bar for the boundaries:
\begin{widetext}
\begin{align}
&\left. \left[\gradrsq D\left(\ne, \frac{d\ne}{dx}\right) \frac{d\ne}{dx} v_e\right] \right|_{x=\xin}^{0} - \int_{\xin}^{0} \gradrsq D\left(\ne, \frac{d\ne}{dx}\right) \frac{d\ne}{dx} v_e' dx = 
\notag \\ & -\int_{\xin}^{0} \ne\Si\bigl(\fsa{\nfc} + \fsa{\ncx}\bigr) v_e  dx
\end{align}
\end{widetext}

We invoke boundary constraints as specified in Sec. \ref{sec:domain}:
\begin{itemize}
\item At the separatrix region boundary ($x=0$), a fixed Dirichlet constraint $\ne(0) = n_{e,0}$ dictates that test functions vanish here ($v_e(0) = 0$), turning the upper boundary term to zero.
\item At the inner domain footprint boundary ($x=\xin$), the behavior is dictated by the configuration flag $\delta_{\mathrm{N}}$ ($\delta_{\mathrm{N}}=1$ for a Neumann scenario, $\delta_{\mathrm{N}}=0$ for a Dirichlet scenario). Under a Neumann setting, the explicitly prescribed flux $\left.\frac{d\ne}{dx}\right|_{\xin} = n'_{e,\mathrm{inner}}$ is substituted into the evaluated boundary term.
\end{itemize}

Multiplying the expression by $-1$, bringing terms onto the left-hand side, and fully expanding the non-linear transport coefficient terms within the volume integral gives the final weak form residual equation. Find $\ne \in V$ such that for all $v_e \in V$:
\begin{widetext}
\begin{align}
&\int_{\xin}^{0} \gradrsq \left[ \frac{C_{\mathrm{ETG}}(x)}{\ne} \frac{d\ne}{dx} + A_{KBM}\frac{dn_e}{dx}+ B_{\mathrm{KBM}}(x) (T_e + T_i) \left(\frac{d\ne}{dx}\right)^2 + B_{\mathrm{KBM}}(x) \ne \frac{d(T_e + T_i)}{dx} \frac{d\ne}{dx} + D_{\mathrm{NEO}}(x) \frac{d\ne}{dx} \right] v_e'  dx 
\notag \\ & - \int_{\xin}^{0} \ne\Si\bigl(\fsa{\nfc}+\fsa{\ncx}\bigr) v_e dx + \delta_{\mathrm{N}}\left. \gradrsq\right|_{x=\xin} \left. D \right|_{x=\xin} n'_{e,\mathrm{inner}} v_e(\xin) = 0
\end{align}
\end{widetext}

Note I am neglecting the neutral pressures here but I think in principle you could include them it might just be an awful differential equation.



\subsection{Non-dimensionalization of Eq.~(A33)}
\label{sec:nondim-A33}

The strong and weak forms derived in App.~\ref{sec:ne-weak-total-pressure}
involve quantities that span many orders of magnitude in SI units
($\ne \!\sim\! 10^{19}$~m$^{-3}$, $\Si \!\sim\! 10^{-14}$~m$^{3}$/s,
$T_{e}\!+\!T_{i} \!\sim\! 10^{-17}$~J,
$D_{\mathrm{NEO}} \!\sim\! 10^{0}$~m$^{2}$/s, etc.). We rescale
Eq.~(A33) onto $\mathcal{O}(1)$ unknowns using a \emph{minimal}
non-dimensionalisation: only four reference scales are chosen freely,
equal to the number of independent SI dimensions in the problem
(length, density, energy, time). All other reference scales follow
from those four by Buckingham's $\pi$-theorem, and the final
dimensionless equation has no remaining multiplicative prefactor.

\subsubsection{Independent reference scales}

The model's four independent SI dimensions are length [m], number
density [m$^{-3}$], energy [J], and time [s]. We pick one numerical
reference constant per dimension, all evaluated at the separatrix:
\begin{subequations}
\label{eq:nondim-scales-A8}
\begin{align}
  L     &\equiv |\xin|
        &&\text{[length, m],} \\
  n_{0} &\equiv n_{e,0}
        &&\text{[density, m$^{-3}$],} \\
  T_{0} &\equiv T_{e}(0) + T_{i}(0)
        &&\text{[energy, J],} \\
  \tau  &\equiv \frac{1}{n_{0}\,S_{0}},\;
              S_{0} \equiv \Si(0)
        &&\text{[time, s].}
\end{align}
\end{subequations}
The choice $\tau = (n_{0}\,S_{0})^{-1}$ is the ionisation time at the
separatrix. It is the only timescale in the model that does not
require a transport coefficient as input, so the resulting
non-dimensionalisation is independent of $D_{\mathrm{NEO}}$,
$C_{\mathrm{ETG}}$, $A_{\mathrm{KBM}}$, or $B_{\mathrm{KBM}}$.

The flux-surface average $\fsa{|\nabla r|^{2}}$ is dimensionless when
$r$ is a length, so it is rescaled by itself,
$\hat{g}(\hat{x}) = \gradrsq$.\footnote{The units listed for
$\fsa{|\nabla r|^{2}}$ in Table~\ref{tab:variables} (m$^{-2}$) are
inconsistent with this convention; the entry should read
``dimensionless''.}
The remaining derived (dependent) scales follow from
\eqref{eq:nondim-scales-A8} algebraically:
\begin{align}
  [V]_{0} &= \frac{L}{\tau} = L\,n_{0}\,S_{0}
          &&\text{[velocity, m/s],}\\
  [D]_{0} &= \frac{L^{2}}{\tau} = L^{2}\,n_{0}\,S_{0}
          &&\text{[diffusivity, m$^{2}$/s],}\\
  [p]_{0} &= n_{0}\,T_{0}
          &&\text{[pressure, Pa].}
\end{align}

\subsubsection{Dimensionless variables}

Define
\begin{align}
  \hat{x}        &= x/L,
  & \hat{n}_{e}    &= \ne/n_{0},
  \\
  \hat{n}_{\mathrm{FC}}  &= \fsa{\nfc}/n_{0},
  & \hat{n}_{\mathrm{CX}} &= \fsa{\ncx}/n_{0},
  \\
  \hat{T}(\hat{x}) &= \frac{T_{e}(x) + T_{i}(x)}{T_{0}},
  & \hat{S}_{i}(\hat{x}) &= \Si(x)/S_{0},
  \\
  \hat{S}_{\mathrm{CX}}(\hat{x}) &= \Scx(x)/S_{0},
  & \hat{g}(\hat{x})     &= \gradrsq.
\end{align}
With $n_{0}=n_{e,0}$, $\hat{x}\in[-1,0]$ and $\hat{n}_{e}(0)=1$
exactly. Differential operators transform as
\begin{align}
  \frac{d}{dx}      &= \frac{1}{L}\frac{d}{d\hat{x}},
  &
  \frac{d\ne}{dx}   &= \frac{n_{0}}{L}\frac{d\hat{n}_{e}}{d\hat{x}}.
  \label{eq:chain-A8}
\end{align}

\subsubsection{Dimensionless transport coefficients}

The four transport quantities in the diffusion definition (Eq.~A29) are
each rescaled so that their natural product in $\hat{D}$ is
dimensionless once divided by $[D]_{0}=L^{2}\,n_{0}\,S_{0}$.

\textit{Neoclassical contribution.}
$D_{\mathrm{NEO}}$ has units of $[D]_{0}$, so directly
\begin{equation}
  \hat{D}_{\mathrm{NEO}}(\hat{x})
  \;=\; \frac{D_{\mathrm{NEO}}(x)}{L^{2}\,n_{0}\,S_{0}}.
  \label{eq:hat-D-NEO}
\end{equation}

\textit{ETG contribution.}
Since $D_{\mathrm{ETG}}=C_{\mathrm{ETG}}/\ne$ has units of $[D]_{0}$,
$C_{\mathrm{ETG}}$ has units of $[D]_{0}\,n_{0}$. Demanding
$D_{\mathrm{ETG}}=[D]_{0}\,\hat{C}_{\mathrm{ETG}}/\hat{n}_{e}$ yields
\begin{equation}
  \hat{C}_{\mathrm{ETG}}(\hat{x})
  \;=\; \frac{C_{\mathrm{ETG}}(x)}{L^{2}\,n_{0}^{2}\,S_{0}}.
  \label{eq:hat-C-ETG}
\end{equation}

\textit{KBM contributions.}
$A_{\mathrm{KBM}}$ has units of $[D]_{0}$, so
\begin{equation}
  \hat{A}_{\mathrm{KBM}}(\hat{x})
  \;=\; \frac{A_{\mathrm{KBM}}(x)}{L^{2}\,n_{0}\,S_{0}}.
  \label{eq:hat-A-KBM}
\end{equation}
$B_{\mathrm{KBM}}$ multiplies $dp/dx$ to give $[D]_{0}$, so its
SI units are
\begin{equation}
  [B_{\mathrm{KBM}}]
  \;=\;\frac{[D]_{0}}{[dp/dx]}
  \;=\;\frac{L^{2}\,n_{0}\,S_{0}}{n_{0}\,T_{0}/L}
  \;=\;\frac{L^{3}\,S_{0}}{T_{0}},
\end{equation}
and therefore
\begin{equation}
  \hat{B}_{\mathrm{KBM}}(\hat{x})
  \;=\; \frac{B_{\mathrm{KBM}}(x)\,T_{0}}{L^{3}\,S_{0}}.
  \label{eq:hat-B-KBM}
\end{equation}

\textit{Term-by-term verification of $\hat{D}$.}
Substituting the dimensionless variables and
Eqs.~\eqref{eq:hat-D-NEO}--\eqref{eq:hat-B-KBM} term-by-term into the
SI definition of $D$ from Eq.~(A29):
\begin{align}
  \frac{C_{\mathrm{ETG}}}{\ne}
  &= \frac{(L^{2}\,n_{0}^{2}\,S_{0})\,\hat{C}_{\mathrm{ETG}}}{n_{0}\,\hat{n}_{e}}
   = (L^{2}\,n_{0}\,S_{0})\,\frac{\hat{C}_{\mathrm{ETG}}}{\hat{n}_{e}},
  \label{eq:term-ETG-A8}
  \\
  A_{\mathrm{KBM}}
  &= (L^{2}\,n_{0}\,S_{0})\,\hat{A}_{\mathrm{KBM}},
  \label{eq:term-A-A8}
  \\
  B_{\mathrm{KBM}}\,(T_{e}{+}T_{i})\,\frac{d\ne}{dx}
  &= \frac{L^{3}\,S_{0}}{T_{0}}\hat{B}_{\mathrm{KBM}}\!\cdot\! T_{0}\hat{T}\!\cdot\!\frac{n_{0}}{L}\frac{d\hat{n}_{e}}{d\hat{x}}
  \notag\\
  &= (L^{2}\,n_{0}\,S_{0})\,\hat{B}_{\mathrm{KBM}}\,\hat{T}\,\frac{d\hat{n}_{e}}{d\hat{x}},
  \label{eq:term-B1-A8}
  \\
  B_{\mathrm{KBM}}\,\ne\,\frac{d(T_{e}{+}T_{i})}{dx}
  &= \frac{L^{3}\,S_{0}}{T_{0}}\hat{B}_{\mathrm{KBM}}\!\cdot\! n_{0}\hat{n}_{e}\!\cdot\!\frac{T_{0}}{L}\frac{d\hat{T}}{d\hat{x}}
  \notag\\
  &= (L^{2}\,n_{0}\,S_{0})\,\hat{B}_{\mathrm{KBM}}\,\hat{n}_{e}\,\frac{d\hat{T}}{d\hat{x}},
  \label{eq:term-B2-A8}
  \\
  D_{\mathrm{NEO}}
  &= (L^{2}\,n_{0}\,S_{0})\,\hat{D}_{\mathrm{NEO}}.
  \label{eq:term-NEO-A8}
\end{align}
Each of the five contributions carries the same dimensional prefactor
$L^{2}\,n_{0}\,S_{0} = [D]_{0}$, so $D = [D]_{0}\,\hat{D}$ with
\begin{widetext}
\begin{align}
  \hat{D}\!\left(\hat{n}_{e},\,\frac{d\hat{n}_{e}}{d\hat{x}}\right)
  \;=\;& \frac{\hat{C}_{\mathrm{ETG}}(\hat{x})}{\hat{n}_{e}}
       + \hat{A}_{\mathrm{KBM}}(\hat{x})
       + \hat{B}_{\mathrm{KBM}}(\hat{x})\,\hat{T}\,\frac{d\hat{n}_{e}}{d\hat{x}}
  \notag \\
       &+ \hat{B}_{\mathrm{KBM}}(\hat{x})\,\hat{n}_{e}\,\frac{d\hat{T}}{d\hat{x}}
       + \hat{D}_{\mathrm{NEO}}(\hat{x}).
  \label{eq:hat-D-A8}
\end{align}
\end{widetext}

\subsubsection{Dimensionless strong form}

Substituting $D=[D]_{0}\,\hat{D}$ together with the chain-rule
relations~\eqref{eq:chain-A8} into the left-hand side of the strong
form (Eq.~A30):
\begin{widetext}
\begin{align}
  \frac{d}{dx}\!\left[\gradrsq\,D\,\frac{d\ne}{dx}\right]
  &= \frac{1}{L}\frac{d}{d\hat{x}}\!\left[\hat{g}\,(L^{2}\,n_{0}\,S_{0})\,\hat{D}\,\frac{n_{0}}{L}\frac{d\hat{n}_{e}}{d\hat{x}}\right]
  \notag \\
  &= \frac{1}{L}\cdot\frac{L^{2}\,n_{0}^{2}\,S_{0}}{L}\,\frac{d}{d\hat{x}}\!\left[\hat{g}\,\hat{D}\,\frac{d\hat{n}_{e}}{d\hat{x}}\right]
  \notag \\
  &= n_{0}^{2}\,S_{0}\,\frac{d}{d\hat{x}}\!\left[\hat{g}\,\hat{D}\,\frac{d\hat{n}_{e}}{d\hat{x}}\right].
  \label{eq:LHS-A8}
\end{align}
\end{widetext}
The right-hand side of Eq.~(A30) becomes
\begin{align}
  -\,\ne\,\Si\,\bigl(\fsa{\nfc}+\fsa{\ncx}\bigr)
  &= -\,(n_{0}\hat{n}_{e})(S_{0}\hat{S}_{i})\,n_{0}\bigl(\hat{n}_{\mathrm{FC}}+\hat{n}_{\mathrm{CX}}\bigr)
  \notag \\
  &= -\,n_{0}^{2}\,S_{0}\,\hat{n}_{e}\,\hat{S}_{i}\,\bigl(\hat{n}_{\mathrm{FC}}+\hat{n}_{\mathrm{CX}}\bigr).
  \label{eq:RHS-A8}
\end{align}
The two factors of $n_{0}^{2}\,S_{0}$ cancel identically, yielding the
dimensionless strong form
\begin{widetext}
\begin{equation}
  \frac{d}{d\hat{x}}\!\left[\hat{g}\,\hat{D}\!\left(\hat{n}_{e},\,\frac{d\hat{n}_{e}}{d\hat{x}}\right)\,\frac{d\hat{n}_{e}}{d\hat{x}}\right]
  \;=\;
  -\,\hat{n}_{e}\,\hat{S}_{i}\,\bigl(\hat{n}_{\mathrm{FC}} + \hat{n}_{\mathrm{CX}}\bigr).
  \label{eq:strong-hat-A8}
\end{equation}
\end{widetext}
There is no residual dimensionless prefactor in
Eq.~\eqref{eq:strong-hat-A8}: the choice $\tau=1/(n_{0}\,S_{0})$ has set
the natural diffusion scale $[D]_{0}=L^{2}/\tau$ \emph{simultaneously}
with the natural source scale, so the diffusive and source contributions
have identical dimensional prefactors.

\subsubsection{Dimensionless weak form}

Multiply Eq.~\eqref{eq:strong-hat-A8} by an arbitrary test function
$v_{e}\in V$ and integrate over $\hat{\Omega} = [-1,0]$:
\begin{widetext}
\begin{equation}
  \int_{-1}^{0} \frac{d}{d\hat{x}}\!\left[\hat{g}\,\hat{D}\,\frac{d\hat{n}_{e}}{d\hat{x}}\right]v_{e}\,d\hat{x}
  \;=\;
  -\int_{-1}^{0} \hat{n}_{e}\,\hat{S}_{i}\,\bigl(\hat{n}_{\mathrm{FC}}+\hat{n}_{\mathrm{CX}}\bigr) v_{e}\,d\hat{x}.
\end{equation}
\end{widetext}
Apply integration by parts to the left-hand side, using
$\int u'\,v\,d\hat{x}=[u\,v]_{-1}^{0}-\int u\,v'\,d\hat{x}$ with
$u=\hat{g}\,\hat{D}\,(d\hat{n}_{e}/d\hat{x})$ and $v=v_{e}$:
\begin{widetext}
\begin{equation}
  \left.\hat{g}\,\hat{D}\,\frac{d\hat{n}_{e}}{d\hat{x}}\,v_{e}\right|_{-1}^{0}
  \;-\;\int_{-1}^{0}\hat{g}\,\hat{D}\,\frac{d\hat{n}_{e}}{d\hat{x}}\,v_{e}'\,d\hat{x}
  \;=\;
  -\int_{-1}^{0}\hat{n}_{e}\,\hat{S}_{i}\,\bigl(\hat{n}_{\mathrm{FC}}+\hat{n}_{\mathrm{CX}}\bigr) v_{e}\,d\hat{x}.
\end{equation}
\end{widetext}
Multiplying through by $-1$ and moving every term to the left-hand side,
\begin{widetext}
\begin{equation}
  \int_{-1}^{0}\hat{g}\,\hat{D}\,\frac{d\hat{n}_{e}}{d\hat{x}}\,v_{e}'\,d\hat{x}
  \;-\;\int_{-1}^{0}\hat{n}_{e}\,\hat{S}_{i}\,\bigl(\hat{n}_{\mathrm{FC}}+\hat{n}_{\mathrm{CX}}\bigr)v_{e}\,d\hat{x}
  \;-\;\left.\hat{g}\,\hat{D}\,\frac{d\hat{n}_{e}}{d\hat{x}}\,v_{e}\right|_{-1}^{0}
  \;=\;0.
  \label{eq:weak-pre-bc-A8}
\end{equation}
\end{widetext}

\textit{Boundary terms.}
At $\hat{x}=0$ the Dirichlet condition $\hat{n}_{e}(0)=1$ forces
$v_{e}(0)=0$, so the upper-end contribution vanishes. At
$\hat{x}=-1$ the inner-boundary contribution is either zero
(Dirichlet, $v_{e}(-1)=0$) or, in the Neumann case, the prescribed
slope is substituted. Defining
\begin{equation}
  \hat{n}'_{e,\mathrm{inner}}
  \;\equiv\; \left.\frac{d\hat{n}_{e}}{d\hat{x}}\right|_{\hat{x}=-1}
  \;=\; \frac{L\,n'_{e,\mathrm{inner}}}{n_{0}},
  \label{eq:hat-flux-A8}
\end{equation}
and $\delta_{\mathrm{N}}=1$ for Neumann (zero for Dirichlet), the
remaining boundary contribution in Eq.~\eqref{eq:weak-pre-bc-A8} is
\begin{equation}
  -\!\left[\hat{g}\,\hat{D}\,\frac{d\hat{n}_{e}}{d\hat{x}}\,v_{e}\right]_{-1}^{0}
  \;=\;
  +\,\delta_{\mathrm{N}}\,\hat{g}(-1)\,\left.\hat{D}\right|_{\hat{x}=-1}\,\hat{n}'_{e,\mathrm{inner}}\,v_{e}(-1).
\end{equation}

\textit{Final residual.}
Substituting the expansion~\eqref{eq:hat-D-A8} of $\hat{D}$ into the
volume integrand of \eqref{eq:weak-pre-bc-A8}, the dimensionless
analogue of Eq.~(A33) is: find $\hat{n}_{e}\in V$ such that for all
$v_{e}\in V$,
\begin{widetext}
\begin{align}
  &\int_{-1}^{0}\!\hat{g}\!\left[
       \frac{\hat{C}_{\mathrm{ETG}}}{\hat{n}_{e}}\frac{d\hat{n}_{e}}{d\hat{x}}
     + \hat{A}_{\mathrm{KBM}}\frac{d\hat{n}_{e}}{d\hat{x}}
     + \hat{B}_{\mathrm{KBM}}\,\hat{T}\!\left(\frac{d\hat{n}_{e}}{d\hat{x}}\right)^{\!2}
     + \hat{B}_{\mathrm{KBM}}\,\hat{n}_{e}\,\frac{d\hat{T}}{d\hat{x}}\,\frac{d\hat{n}_{e}}{d\hat{x}}
     + \hat{D}_{\mathrm{NEO}}\,\frac{d\hat{n}_{e}}{d\hat{x}}
  \right]\!v_{e}'\,d\hat{x}
  \notag \\
  & \quad
  -\,\int_{-1}^{0}\,\hat{n}_{e}\,\hat{S}_{i}\,\bigl(\hat{n}_{\mathrm{FC}}+\hat{n}_{\mathrm{CX}}\bigr)v_{e}\,d\hat{x}
  \;+\;\delta_{\mathrm{N}}\,\hat{g}(-1)\,\left.\hat{D}\right|_{\hat{x}=-1}\,\hat{n}'_{e,\mathrm{inner}}\,v_{e}(-1)
  \;=\;0.
  \label{eq:weak-hat-A8}
\end{align}
\end{widetext}
No multiplicative dimensionless parameter appears anywhere in
Eq.~\eqref{eq:weak-hat-A8}; the only dimensionless \emph{coefficients}
are $\hat{C}_{\mathrm{ETG}}$, $\hat{A}_{\mathrm{KBM}}$,
$\hat{B}_{\mathrm{KBM}}$, $\hat{D}_{\mathrm{NEO}}$, each carrying
genuine physics information.

\subsubsection{Dimensionless boundary conditions}

Translating each Dirichlet/Neumann condition from
App.~\ref{sec:domain} into the rescaled variables:
\begin{align}
  \hat{n}_{e}(0)        &= 1, \\
  \hat{n}_{\mathrm{FC}}(0) &= \fsa{\nfc}_{0}/n_{0}, \\
  \hat{n}_{\mathrm{CX}}(0) &= \fsa{\ncx}_{0}/n_{0},
\end{align}
and at the inner boundary,
\begin{equation}
  \hat{n}_{e}(-1) \;=\; n_{e,\mathrm{inner}}/n_{0}
  \quad\text{(Dirichlet)}
\end{equation}
or Eq.~\eqref{eq:hat-flux-A8} (Neumann).

\subsubsection{Order-of-magnitude estimates}

For typical DIII-D pedestal parameters,
$L \!\sim\! 5\!\times\!10^{-2}$~m,
$n_{0} \!\sim\! 5\!\times\!10^{19}$~m$^{-3}$,
$T_{0} \!\sim\! 200$~eV $\!\approx\! 3.2\!\times\!10^{-17}$~J, and
$S_{0} \!\sim\! 3\!\times\!10^{-14}$~m$^{3}$/s. The derived scales
follow:
\begin{align}
  \tau    &= \frac{1}{n_{0}\,S_{0}}
            \;=\; \frac{1}{(5\!\times\!10^{19})(3\!\times\!10^{-14})}
            \;\approx\; 6.7\!\times\!10^{-7}~\text{s},
  \\
  [V]_{0} &= \frac{L}{\tau}
            \;=\; \frac{5\!\times\!10^{-2}}{6.7\!\times\!10^{-7}}
            \;\approx\; 7.5\!\times\!10^{4}~\text{m/s},
  \\
  [D]_{0} &= \frac{L^{2}}{\tau}
            \;=\; \frac{(5\!\times\!10^{-2})^{2}}{6.7\!\times\!10^{-7}}
            \;\approx\; 3.8\!\times\!10^{3}~\text{m}^{2}/\text{s}.
\end{align}
With representative SI transport coefficients
$D_{\mathrm{NEO}} \!\sim\! 0.5$~m$^{2}$/s and
$C_{\mathrm{ETG}} \!\sim\! D_{\mathrm{ETG}}\,\ne \!\sim\! 2.5\!\times\!10^{19}$~m$^{-1}$\,s$^{-1}$,
the dimensionless coefficients are
\begin{align}
  \hat{D}_{\mathrm{NEO}}
  &\;\approx\; \frac{0.5}{3.8\!\times\!10^{3}}
   \;\approx\; 1.3\!\times\!10^{-4},
  \\
  \hat{C}_{\mathrm{ETG}}
  &\;\approx\; \frac{2.5\!\times\!10^{19}}{(2.5\!\times\!10^{-3})(2.5\!\times\!10^{39})(3\!\times\!10^{-14})}
  \notag\\
  &\;\approx\; 1.3\!\times\!10^{-4}.
\end{align}
The smallness of the dimensionless diffusion contributions reflects
the physically meaningful fact that the pedestal is a thin boundary
layer in which slow diffusion competes with fast ionisation. The
high-derivative term in Eq.~\eqref{eq:strong-hat-A8} carries a small
prefactor but is balanced by large $d\hat{n}_{e}/d\hat{x}$ inside the
boundary layer (singular-perturbation behaviour); the equation itself
is well-posed without further rescaling.

\subsubsection{Optional numerical preconditioning}
\label{sec:precond-A8}

For the SNES/Newton solves of Eq.~\eqref{eq:weak-hat-A8} it is sometimes
useful to introduce an additional, purely numerical, scale $D_{0}$ that
brings each assembled coefficient near unity in floating point.
Defining
\begin{align}
  \tilde{D}_{\mathrm{NEO}} &= \frac{D_{\mathrm{NEO}}}{D_{0}}, &
  \tilde{C}_{\mathrm{ETG}} &= \frac{C_{\mathrm{ETG}}}{D_{0}\,n_{0}},
  \\
  \tilde{A}_{\mathrm{KBM}} &= \frac{A_{\mathrm{KBM}}}{D_{0}}, &
  \tilde{B}_{\mathrm{KBM}} &= \frac{B_{\mathrm{KBM}}\,n_{0}\,T_{0}}{L\,D_{0}},
\end{align}
the strong form acquires a single Damk\"ohler-like prefactor:
\begin{equation}
  \frac{d}{d\hat{x}}\!\left[\hat{g}\,\tilde{D}\,\frac{d\hat{n}_{e}}{d\hat{x}}\right]
  \;=\;-\,\Lambda\,\hat{n}_{e}\,\hat{S}_{i}\,(\hat{n}_{\mathrm{FC}}+\hat{n}_{\mathrm{CX}}),
\end{equation}
\begin{equation}
  \Lambda \;\equiv\; \frac{n_{0}\,S_{0}\,L^{2}}{D_{0}}
  \;=\; \frac{[D]_{0}}{D_{0}}.
\end{equation}
Choosing $D_{0}=D_{\mathrm{NEO}}(0)\sim 0.5$~m$^{2}$/s gives
$\Lambda \sim 7.5\!\times\!10^{3}$ but makes
$\tilde{D}_{\mathrm{NEO}}\sim\tilde{D}_{\mathrm{ETG}}\sim 1$, which can
be preferable for the floating-point Jacobian. Setting
$D_{0}=[D]_{0}=L^{2}/\tau$ recovers the foundational form
\eqref{eq:strong-hat-A8} ($\Lambda=1$). Both choices yield identical
SI solutions; the difference is only in the conditioning of the linear
solves.

\subsubsection{Non-dimensionalisation of the neutral equations}
\label{sec:nondim-neutrals}

The same four reference scales \eqref{eq:nondim-scales-A8} are used to
rescale the FC- and CX-neutral fluid equations of
App.~\ref{sec:domain}.  No new scales are introduced.

\paragraph{Velocity-scale rescalings.}
The Franck--Condon streaming speed
$|V_{\mathrm{FC}}| = \sqrt{8\,E_{\mathrm{FC}}/(\pi^{2}\,M_{\mathrm{eff}}\,M_{i})}$
is a constant of the equilibrium, while the CX speed
$|V_{\mathrm{CX}}|(x)$ is the local ion thermal speed.  Both have units
of velocity, so we rescale by $[V]_{0} = L/\tau = L\,n_{0}\,S_{0}$:
\begin{align}
  \hat V_{\mathrm{FC}} &\equiv \frac{|V_{\mathrm{FC}}|}{L\,n_{0}\,S_{0}}, &
  \hat V_{\mathrm{CX}}(\hat x) &\equiv \frac{|V_{\mathrm{CX}}|(x)}{L\,n_{0}\,S_{0}}.
  \label{eq:hat-V-A8}
\end{align}
The form factors $f_{\mathrm{FC}}$ and $f_{\mathrm{CX}}$ are already
dimensionless and are unchanged.

\paragraph{FC-neutral equation.}
Starting from the SI strong form (cf.\ App.~\ref{sec:domain})
\begin{equation}
  |V_{\mathrm{FC}}|\,\frac{d}{dx}\!\bigl[f_{\mathrm{FC}}\,\fsa{|\nabla r|^{2}}\,\fsa{\nfc}\bigr]
  \;=\;\ne\,(\Si + \Scx)\,\fsa{\nfc},
\end{equation}
substitute $|V_{\mathrm{FC}}| = (L\,n_{0}\,S_{0})\,\hat V_{\mathrm{FC}}$,
$d/dx = L^{-1}\,d/d\hat x$, $\fsa{\nfc} = n_{0}\,\hat n_{\mathrm{FC}}$,
$\hat g = \fsa{|\nabla r|^{2}}$, $\ne = n_{0}\,\hat n_{e}$,
$\Si = S_{0}\,\hat S_{i}$, $\Scx = S_{0}\,\hat S_{\mathrm{CX}}$:
\begin{align}
  \mathrm{LHS}
  &= (L\,n_{0}\,S_{0})\,\hat V_{\mathrm{FC}}
     \cdot\frac{1}{L}\,\frac{d}{d\hat x}\!\bigl[f_{\mathrm{FC}}\,\hat g\,n_{0}\,\hat n_{\mathrm{FC}}\bigr]
   \notag\\
  &= n_{0}^{2}\,S_{0}\,\hat V_{\mathrm{FC}}\,\frac{d}{d\hat x}\!\bigl[f_{\mathrm{FC}}\,\hat g\,\hat n_{\mathrm{FC}}\bigr],\\
  \mathrm{RHS}
  &= (n_{0}\,\hat n_{e})(S_{0}\,\hat S_{i} + S_{0}\,\hat S_{\mathrm{CX}})(n_{0}\,\hat n_{\mathrm{FC}})
   \notag\\
  &= n_{0}^{2}\,S_{0}\,\hat n_{e}\,(\hat S_{i}+\hat S_{\mathrm{CX}})\,\hat n_{\mathrm{FC}}.
\end{align}
The common factor $n_{0}^{2}\,S_{0}$ cancels, leaving the dimensionless
strong form
\begin{equation}
  \hat V_{\mathrm{FC}}\,\frac{d}{d\hat x}\!\bigl[f_{\mathrm{FC}}\,\hat g\,\hat n_{\mathrm{FC}}\bigr]
  \;=\; \hat n_{e}\,(\hat S_{i}+\hat S_{\mathrm{CX}})\,\hat n_{\mathrm{FC}}.
  \label{eq:strong-hat-nFC-A8}
\end{equation}

\paragraph{CX-neutral equation.}
The SI strong form (cf.\ App.~\ref{sec:domain})
\begin{equation}
  \frac{d}{dx}\!\bigl[|V_{\mathrm{CX}}|\,f_{\mathrm{CX}}\,\fsa{|\nabla r|^{2}}\,\fsa{\ncx}\bigr]
  \;=\;\ne\!\left(\Si\,\fsa{\ncx} - \tfrac{1}{2}\,\Scx\,\fsa{\nfc}\right)
\end{equation}
rescales analogously.  Substituting and combining the
$\frac{1}{L}\!\cdot\! L$ from $d/dx$ and $|V_{\mathrm{CX}}|$,
\begin{align}
  \mathrm{LHS}
  &= \frac{1}{L}\,\frac{d}{d\hat x}\!\bigl[(L\,n_{0}\,S_{0})\hat V_{\mathrm{CX}}\,f_{\mathrm{CX}}\,\hat g\,n_{0}\,\hat n_{\mathrm{CX}}\bigr]
   \notag\\
  &= n_{0}^{2}\,S_{0}\,\frac{d}{d\hat x}\!\bigl[\hat V_{\mathrm{CX}}\,f_{\mathrm{CX}}\,\hat g\,\hat n_{\mathrm{CX}}\bigr],\\
  \mathrm{RHS}
  &= n_{0}^{2}\,S_{0}\,\hat n_{e}\!\left(\hat S_{i}\,\hat n_{\mathrm{CX}}
       - \tfrac{1}{2}\,\hat S_{\mathrm{CX}}\,\hat n_{\mathrm{FC}}\right).
\end{align}
The common $n_{0}^{2}\,S_{0}$ cancels, giving
\begin{equation}
  \frac{d}{d\hat x}\!\bigl[\hat V_{\mathrm{CX}}\,f_{\mathrm{CX}}\,\hat g\,\hat n_{\mathrm{CX}}\bigr]
  \;=\; \hat n_{e}\!\left(\hat S_{i}\,\hat n_{\mathrm{CX}}
       - \tfrac{1}{2}\,\hat S_{\mathrm{CX}}\,\hat n_{\mathrm{FC}}\right).
  \label{eq:strong-hat-nCX-A8}
\end{equation}

\paragraph{Weak forms (matching the strong-residual convention used in
\texttt{solver.py}).}
Multiplying \eqref{eq:strong-hat-nFC-A8} by a test function
$v_{\mathrm{F}}\in V$ and integrating over $\hat\Omega = [-1,0]$
(no integration by parts) gives
\begin{equation}
  \int_{-1}^{0}\!\!\left[\hat V_{\mathrm{FC}}\,\frac{d}{d\hat x}\!\bigl[f_{\mathrm{FC}}\,\hat g\,\hat n_{\mathrm{FC}}\bigr]
       - \hat n_{e}\,(\hat S_{i}+\hat S_{\mathrm{CX}})\,\hat n_{\mathrm{FC}}\right]\!v_{\mathrm{F}}\,d\hat x
  \;=\; 0.
  \label{eq:weak-hat-nFC-A8}
\end{equation}
Similarly, with $v_{\mathrm{C}}\in V$ for \eqref{eq:strong-hat-nCX-A8},
\begin{equation}
  \int_{-1}^{0}\!\!\left[\frac{d}{d\hat x}\!\bigl[\hat V_{\mathrm{CX}}\,f_{\mathrm{CX}}\,\hat g\,\hat n_{\mathrm{CX}}\bigr]
       - \hat n_{e}\!\left(\hat S_{i}\,\hat n_{\mathrm{CX}}
              - \tfrac{1}{2}\,\hat S_{\mathrm{CX}}\,\hat n_{\mathrm{FC}}\right)\right]\!v_{\mathrm{C}}\,d\hat x
  \;=\; 0.
  \label{eq:weak-hat-nCX-A8}
\end{equation}
Both equations are first-order, with characteristics flowing from the
separatrix inward, so the single Dirichlet inflow conditions
\begin{align}
  \hat n_{\mathrm{FC}}(0) &= \fsa{\nfc}_{0}/n_{0}, &
  \hat n_{\mathrm{CX}}(0) &= \fsa{\ncx}_{0}/n_{0},
\end{align}
suffice to close the system together with the two $\hat n_{e}$
boundary conditions of Sec.~\ref{sec:nondim-A33}.

\subsubsection{Implementation roadmap}

Implementing Eqs.~\eqref{eq:weak-hat-A8},
\eqref{eq:weak-hat-nFC-A8}--\eqref{eq:weak-hat-nCX-A8} in
\texttt{solver-nondim.py} requires:
\begin{enumerate}
  \item Rebuilding the Firedrake \texttt{IntervalMesh} on $[-1,0]$
    (or, equivalently, storing $L$ and dividing every gradient by $L$
    inside the residual).
  \item Replacing each \texttt{Function} stored in \texttt{\_fd\_cache}
    (\texttt{C\_ETG\_fd}, \texttt{D\_NEO\_fd}, \texttt{A\_KBM\_fd},
    \texttt{B\_KBM\_fd}, \texttt{T\_tot\_fd}, \texttt{dT\_tot\_dx\_fd},
    \texttt{Si\_fd}, \texttt{Scx\_fd}, \texttt{Vcx\_fd}, \texttt{g\_fd})
    by its rescaled counterpart per
    Eqs.~\eqref{eq:hat-D-NEO}--\eqref{eq:hat-B-KBM} and
    \eqref{eq:hat-V-A8}.  In particular,
    $\hat S_{i,\mathrm{CX}} = S_{i,\mathrm{CX}}/S_{0}$,
    $\hat T = (T_{e}+T_{i})/T_{0}$, and $\hat g = \fsa{|\nabla r|^{2}}$
    are dimensionless and can be evaluated directly from their SI
    inputs.
  \item Replacing the Dirichlet/Neumann boundary data with
    $\hat n_{e}(0)=1$,
    $\hat n_{e}(-1)=n_{e,\mathrm{inner}}/n_{0}$,
    $\hat n'_{e,\mathrm{inner}}=L\,n'_{e,\mathrm{inner}}/n_{0}$,
    $\hat n_{\mathrm{FC}}(0)=\fsa{\nfc}_{0}/n_{0}$, and
    $\hat n_{\mathrm{CX}}(0)=\fsa{\ncx}_{0}/n_{0}$.
  \item Assembling the three residuals
    \eqref{eq:weak-hat-A8},
    \eqref{eq:weak-hat-nFC-A8}--\eqref{eq:weak-hat-nCX-A8} into
    a single mixed Newton residual $\hat F(\hat u) = 0$ on
    $\hat W = \hat V \times \hat V \times \hat V$ and solving with
    Firedrake's SNES driver.
  \item Recovering the SI solution at the end of
    \texttt{solve\_coupled\_nondim} via
    $\ne(x)=n_{0}\,\hat n_{e}(x/L)$,
    $\fsa{\nfc}(x)=n_{0}\,\hat n_{\mathrm{FC}}(x/L)$,
    $\fsa{\ncx}(x)=n_{0}\,\hat n_{\mathrm{CX}}(x/L)$.
\end{enumerate}
No SNES source-term scaling is required (since no $\Lambda$ appears in
Eqs.~\eqref{eq:weak-hat-A8}, \eqref{eq:weak-hat-nFC-A8}--\eqref{eq:weak-hat-nCX-A8}).
If the floating-point conditioning of the foundational form proves
insufficient, the preconditioning of Sec.~\ref{sec:precond-A8} can be
applied to the $\hat n_{e}$ diffusion coefficients only, leaving the
source integrand and the neutral equations as written.



\section{Variable identities}
\label{app:Variables}

Tables~\ref{tab:variables} and~\ref{tab:variables-cont} list symbols and parameters appearing in the main text and appendices.
Flux-surface-averaged quantities use $\langle\cdot\rangle$ notation as in \cite{saarelma2023testing}.

\begin{table*}[t]
\caption{\label{tab:variables}Symbols and parameters used in the Saarelma--Connor pedestal study (part~I).}
\begin{ruledtabular}
\begin{tabular}{lll}
Symbol & Description & Units \\
\hline
\multicolumn{3}{l}{\textit{Coordinates and geometry}} \\
$x$ & Radial coordinate measured inward from separatrix ($x=0$ at separatrix, $x<0$ inwards). & m \\
$r$ & Minor radius. & m \\
$r_{\mathrm{sep}}$ & Minor radius of the separatrix. & m \\
$x_{\mathrm{inner}}$ & Inner edge of the 1D pedestal domain (pedestal foot). & m \\
$\psi_N$ & Normalized poloidal flux, used to label flux surfaces. & dimensionless \\
$\theta$ & Poloidal angle in flux coordinates. & rad \\
$R$ & Major radius. & m \\
$a$ & Device minor radius used in pedestal transport scalings. & m \\
$\Omega$ & Computational interval $[x_{\mathrm{inner}},\,0]$. & m \\
$\langle A\rangle$ & Flux-surface average, $\oint R^2 A\,d\theta / \oint R^2\,d\theta$. & same as $A$ \\
$\langle|\nabla r|^2\rangle$ & Flux-surface average of the metric factor $|\nabla r|^2$. & m$^{-2}$ \\
\hline
\multicolumn{3}{l}{\textit{Dependent fields (three-fluid model)}} \\
$n_e$ & Electron number density. & m$^{-3}$ \\
$n_{\mathrm{FC}}$ & Franck--Condon (FC) neutral number density. & m$^{-3}$ \\
$n_{\mathrm{CX}}$ & Charge-exchange (CX) neutral number density. & m$^{-3}$ \\
$\langle n_{\mathrm{FC}}\rangle$ & Flux-surface-averaged FC neutral density solved in the coupled model. & m$^{-3}$ \\
$\langle n_{\mathrm{CX}}\rangle$ & Flux-surface-averaged CX neutral density solved in the coupled model. & m$^{-3}$ \\
$n_{e,\mathrm{pres}}$ & Experimental electron density from the p-file (input profile). & m$^{-3}$ \\
\hline
\multicolumn{3}{l}{\textit{Transport coefficients and reaction rates}} \\
$D_{\mathrm{ped}}$ & Pedestal electron diffusion coefficient ($D_{\mathrm{KBM}}+D_{\mathrm{ETG}}+D_{\mathrm{NEO}}$). & m$^2$\,s$^{-1}$ \\
$D_{\mathrm{KBM}}$ & Kinetic-ballooning-mode contribution to $D_{\mathrm{ped}}$. & m$^2$\,s$^{-1}$ \\
$D_{\mathrm{ETG}}$ & Electron-temperature-gradient contribution to $D_{\mathrm{ped}}$. & m$^2$\,s$^{-1}$ \\
$D_{\mathrm{NEO}}$ & Neoclassical contribution to $D_{\mathrm{ped}}$. & m$^2$\,s$^{-1}$ \\
$S_i$ & Electron-impact ionization rate coefficient. & m$^3$\,s$^{-1}$ \\
$S_{\mathrm{CX}}$ & Charge-exchange rate coefficient. & m$^3$\,s$^{-1}$ \\
$|V_{\mathrm{FC}}|$ & Magnitude of the FC neutral streaming speed. & m\,s$^{-1}$ \\
$|V_{\mathrm{CX}}|$ & Magnitude of the CX neutral streaming speed (local ion thermal speed). & m\,s$^{-1}$ \\
$f_{\mathrm{FC}}$ & FC form factor relating poloidal structure to $\langle n_{\mathrm{FC}}\rangle$ (unity in this study). & dimensionless \\
$f_{\mathrm{CX}}$ & CX form factor relating poloidal structure to $\langle n_{\mathrm{CX}}\rangle$ (unity in this study). & dimensionless \\
$\alpha$ & Normalized pressure-gradient parameter controlling KBM onset. & dimensionless \\
$c_s$ & Ion sound speed on flux surfaces. & m\,s$^{-1}$ \\
$\rho_s$ & Ion gyroradius on flux surfaces. & m \\
\hline
\multicolumn{3}{l}{\textit{Boundary conditions}} \\
$n_{e,0}$ & Electron density at the separatrix ($x=0$). & m$^{-3}$ \\
$\langle n_{\mathrm{FC}}\rangle_0$ & FC neutral density at the separatrix. & m$^{-3}$ \\
$\langle n_{\mathrm{CX}}\rangle_0$ & CX neutral density at the separatrix. & m$^{-3}$ \\
$n_{e,\mathrm{inner}}$ & Prescribed electron density at $x_{\mathrm{inner}}$ (Dirichlet option). & m$^{-3}$ \\
$\left.{dn_e}/{dx}\right|_{x_{\mathrm{inner}}}$ & Prescribed electron density gradient at $x_{\mathrm{inner}}$ (Neumann option). & m$^{-4}$ \\
$\psi_{N,\mathrm{inner}}$ & Normalized flux label where the inner boundary is placed. & dimensionless \\
\hline
\multicolumn{3}{l}{\textit{Model free parameters (this study)}} \\
$P_{\mathrm{tot},e}$ & Electron heating power used in the ETG diffusion term. & W \\
$\alpha_{\mathrm{crit}}$ & Critical $\alpha$ above which $D_{\mathrm{KBM}}$ is activated. & dimensionless \\
$C_{\mathrm{KBM}}$ & KBM diffusion scaling coefficient. & m$^2$\,s$^{-1}$ \\
$D_{e,\chi i,\mathrm{ETG}}$ & ETG diffusion tuning parameter (\texttt{De\_chie\_etg} in code). & m$^2$\,s$^{-1}$ \\
$n_{\mathrm{FC},0}$ & Input FC neutral density at the separatrix (\texttt{nFC\_x0}). & m$^{-3}$ \\
\end{tabular}
\end{ruledtabular}
\end{table*}

\begin{table*}[t]
\caption{\label{tab:variables-cont}Symbols and parameters used in the Saarelma--Connor pedestal study (part~II).}
\begin{ruledtabular}
\begin{tabular}{lll}
Symbol & Description & Units \\
\hline
\multicolumn{3}{l}{\textit{Physical constants and species data}} \\
$E_{\mathrm{FC}}$ & Franck--Condon dissociation energy of molecular deuterium ($\sim 3$\,eV). & J \\
$Z_i$ & Main-ion charge number. & dimensionless \\
$M_i$ & Proton (nucleon) mass. & kg \\
$M_{\mathrm{eff}}$ & Effective ion mass for D or D--T mixtures. & kg \\
$M_e$ & Electron mass. & kg \\
\hline
\multicolumn{3}{l}{\textit{Numerical discretization and solvers}} \\
$p$ & Polynomial degree of the continuous Galerkin finite-element space. & dimensionless \\
\texttt{mesh\_n} & Number of cells on the 1D Firedrake mesh. & dimensionless \\
\texttt{x\_res} & Number of radial points for precomputed coefficient grids. & dimensionless \\
$\lambda$ & Picard under-relaxation factor (\texttt{relax}). & dimensionless \\
\texttt{picard\_tol} & Relative convergence tolerance on $n_e$ between Picard iterates. & dimensionless \\
\texttt{max\_picard} & Maximum number of Picard iterations. & dimensionless \\
$v_e,\,v_F,\,v_C$ & Finite-element test functions for $n_e$, $\langle n_{\mathrm{FC}}\rangle$, $\langle n_{\mathrm{CX}}\rangle$. & dimensionless \\
$F_1,\,F_2,\,F_3$ & Variational residuals for Eqs.~(8)--(10) in the Firedrake implementation. & N/A \\
\hline
\multicolumn{3}{l}{\textit{Reduced-model auxiliary quantity (original solver only)}} \\
$L_2(n_e)$ & Integral operator appearing in the collapsed Saarelma--Connor equation before full integration. & N/A \\
\end{tabular}
\end{ruledtabular}
\end{table*}

\end{document}